\documentclass[letterpaper, 12pt]{report}
\usepackage[utf8]{inputenc}
\usepackage[english, spanish]{babel}
\usepackage{fullpage} % changes the margin
\usepackage{graphicx} 
\usepackage{enumitem} 
\usepackage{chngcntr}
\usepackage{multirow}
\counterwithin{figure}{section}
\renewcommand{\thesection}{\arabic{section}} 
\renewcommand{\thesubsection}{\thesection.\arabic{subsection}}
\renewcommand{\baselinestretch}{1.5}
\usepackage{float}
\bibliographystyle{apalike}
\setlength\belowcaptionskip{10pt}
\linespread{1.5}

\raggedbottom{}

\begin{document}

\begin{titlepage}
	\centering
	\includegraphics[width=0.3\textwidth]{Images/logo_utb.png}\par\vspace{1cm}
	{\scshape\LARGE Universidad Tecnológica de Bolívar \par}
	\vspace{1cm}

	{\scshape\Large FÍSICA ELÉCTRICA \par}
	\vspace{.2cm}

	% chktex-file 8
	{\scshape\Large H1 - C \par}
	\vspace{1cm}
	% chktex-file 8
	\slshape {\Large \bfseries{}Informe de Laboratorio No. 2 \\}
	\vspace{1cm}

	\slshape {\itshape{} Mauro González, T00067622 \\}
	\slshape {\itshape{} German De Armas Castaño, T00068765 \\}
	\slshape {\itshape{} Angel Vega Rodriguez, T00068186 \\}
	\slshape {\itshape{} Juan Jose Osorio Ariza, T00067316 \\}
	\slshape {\itshape{} Juan Eduardo barón, T00065901 \\}
	\vfill
	Revisado Por \\
	Gabriel Hoyos Gomez Casseres\\
	{\large \today\par}
\end{titlepage}

% ----------------------------------------------------------------------|>
\section{Introducción}

En el siguiente informe presentamos los diferentes métodos implementados en
la práctica de laboratorio consolidados en la guía, los cuales nos permitieron
de forma práctica obtener mediciones de corriente, potencial y resistencia a
través del multímetro (instrumento encargado de realizar dichas mediciones).

\vspace{.5cm}

Gracias a los resultados obtenidos se observa la diferencia en las medidas
tomadas con el multímetro digital y el análogo, los cuales tienen diferente
exactitud. A su vez, la experiencia en el laboratorio permite manejar con
mejor claridad el funcionamiento de los circuitos y la importancia que tienen
los capacitores.

\newpage

% ----------------------------------------------------------------------|>
\section{Objetivos}

\subsection{Objetivos General}

\begin{itemize}
	\item Reconocer los diferentes elementos en el panel Fuentes y aprender
	      cómo Usarlos
	\item Aprender a medir voltaje, corriente y resistencia con un multímetro
	      (analógico y digital).
\end{itemize}

\subsection{Objetivos específicos}

\begin{itemize}
	\item Reconocer la diferencia entre corriente y resistencia.
	\item Diferenciar los diferentes usos  del voltímetro, ohmímetro y
	      amperímetro
\end{itemize}

\newpage

% ----------------------------------------------------------------------|>
\section{Marco Teórico}

\subsection{Circuito eléctrico}

se considera un circuito eléctrico al flujo de cargas siguen una trayectoria
al desplazarse o, en otras palabras, es el recorrido estipula que tendrá una
corriente eléctrica.

\subsection{Diferencia de potencial}

Es una ley básica de los circuitos eléctricos que
establece la diferencia de potencial V que aplicamos entre los extremos de
un conductor determinado es proporcional a la intensidad de la corriente I
que circula por el citado conductor [2]. Ohm completó la ley introduciendo
la noción de resistencia eléctrica R, que es el factor de proporcionalidad
que aparece en la relación entre V y I:\@
V $\to$ R * I

\vspace{.5cm}

También se puede definir como el trabajo por unidad de carga ejercido por el
campo eléctrico sobre una partícula cargada para moverla entre dos posiciones
determinadas.

\subsection{Voltaje}

También llamada como la diferencia de potencial eléctrico y representada por
la letra ``V”, es el trabajo que se debe emplear en una cantidad de carga para
que esta logre desplazarse o impulsarse de un punto a otro.

\subsection{Fuente de corriente alterna}

: Es un tipo de corriente eléctrica, en la que la dirección del flujo
de electrones va y viene a intervalos regulares o en ciclos. La corriente
alterna es generada en las grandes centrales eléctricas por medio de
alternadores, los que a su vez alimentan a los transformadores de transmisión.
~\cite{CorrienteAlterna/Continua}

\subsection{Fuente de corriente continua}

Es la corriente eléctrica que fluye de forma constante en una dirección,
como la que fluye en una linterna o en cualquier otro aparato con baterías
es corriente continua.~\cite{CorrienteAlterna/Continua}

\subsection{Resistencia}

Se define la resistencia eléctrica como la mayor o menor dificultad que
opone un cuerpo al de la corriente eléctrica. Los materiales que presentan
una gran oposición al paso de la electricidad reciben el nombre de
aislante, y, en consecuencia, tienen una elevada resistencia eléctrica.

\smallskip{}

Por el contrario, se llaman conductores a los materiales que apenas
oponen resistencia al paso de la corriente

\subsection{Amperímetro}

Es un instrumento que sirve para medir la intensidad de la corriente 
eléctrica al conectarse a un circuito en serie.

\subsection{Voltímetro}

Mide la diferencia de potencial entre dos puntos de un circuito eléctrico 
y por lo tanto, se debe conectar en paralelo con la porción del circuito 
sobre el que se quiere realizar la medida. 



% ------------------------------------------------ |>
\subsection{Fórmula de la desviación estándar}

\includegraphics[scale = .6]{./Images/FormulaDesviacionEstandar.png} \hfill
\break{}~\cite{ImagenDesviacionEstandar}
% chktex-file 24
\label{img:desviacionEstandar}


% ----------------------------------------------------------------------|>
% \section{Montaje Experimental}

% ----------------------------------------------------------------------|>
\section{Datos Experimentales}

\begin{table}[h]
	\begin{center}
		% chktex-file 44
		\begin{tabular}{|c|c|c|}
			\multicolumn{3}{c}{Resistencias}              \\ \hline
			Medición & Código        & Multímetro digital \\ \hline
			1        & 220 $\pm$ 5\% & 217,9 $\Omega$     \\ \hline
			2        & 150 $\pm$ 5\% & 149,1 $\Omega$     \\ \hline
			3        & 470 $\pm$ 5\% & 465,0 $\Omega$     \\ \hline
			4        & 330 $\pm$ 5\% & 324,3 $\Omega$     \\ \hline
		\end{tabular}
		\caption{Tabla de resistencias}
		% chktex-file 24
		\label{tab:resistencias}
	\end{center}
\end{table}

% ---------- TABLAS DE VOLTAJE EN AC ---------- |>

\begin{table}[H]
	\begin{center}
		% chktex-file 44
		\begin{tabular}{|c|c|}
			\multicolumn{2}{c}{Voltaje en AC}                     \\ \hline
			Multímetro digital & Multímetro analógico (1/Persona) \\ \hline
			5,44 V             & 6,80 V                           \\
			                   & 6,70 V                           \\
			                   & 6,70 V                           \\
			                   & 6,80 V                           \\
			                   & 6,60 V                           \\ \hline
			Promedio           & 6,72 V                           \\ \hline
		\end{tabular}
		\caption{Voltaje en AC [Primera Medición]}
		% chktex-file 24
		\label{tab:voltajeAC_1}
	\end{center}
\end{table}

\begin{table}[H]
	\begin{center}
		% chktex-file 44
		\begin{tabular}{|c|c|}
			\multicolumn{2}{c}{Voltaje en AC}                     \\ \hline
			Multímetro digital & Multímetro analógico (1/Persona) \\ \hline
			17,17 V            & 21,08 V                          \\
			                   & 20,75 V                          \\
			                   & 21,75 V                          \\
			                   & 20,05 V                          \\
			                   & 21,10 V                          \\ \hline
			Promedio           & 20,95 V                          \\ \hline
		\end{tabular}
		\caption{Voltaje en AC [Segunda Medición]}
		% chktex-file 24
		\label{tab:voltajeAC_2}
	\end{center}
\end{table}

% ---------- TABLAS DE VOLTAJE EN DC ---------- |>

\begin{table}[H]
	\begin{center}
		% chktex-file 44
		\begin{tabular}{|c|c|}
			\multicolumn{2}{c}{Voltaje en DC}                     \\ \hline
			Multímetro digital & Multímetro analógico (1/Persona) \\ \hline
			6,88 V             & 8,20 V                           \\
			                   & 8,20 V                           \\
			                   & 8,20 V                           \\
			                   & 8,20 V                           \\
			                   & 8,15 V                           \\ \hline
			Promedio           & 8.19 V                           \\ \hline
		\end{tabular}
		\caption{Voltaje en DC [Primera Medición]}
		% chktex-file 24
		\label{tab:voltajeDC_1}
	\end{center}
\end{table}

\begin{table}[H]
	\begin{center}
		% chktex-file 44
		\begin{tabular}{|c|c|}
			\multicolumn{2}{c}{Voltaje en DC}                     \\ \hline
			Multímetro digital & Multímetro analógico (1/Persona) \\ \hline
			17,47 V            & 18,85 V                          \\
			                   & 18,50 V                          \\
			                   & 18,50 V                          \\
			                   & 19,00 V                          \\
			                   & 18,50 V                          \\ \hline
			Promedio           & 18,67 V                          \\ \hline
		\end{tabular}
		\caption{Voltaje en DC [Segunda Medición]}
		% chktex-file 24
		\label{tab:voltajeDC_2}
	\end{center}
\end{table}

% ----------------------------------------------------------------------|>
\section{Análisis de datos}

\subsection{Error en las mediciones de voltaje AC y DC}

Usando la formula de la desviación estándar,~\ref{img:desviacionEstandar},
se obtiene:

\begin{table}[H]
	\begin{center}
		% chktex-file 44
		\begin{tabular}{|c|c|}
			\multicolumn{2}{c}{Voltaje en AC}               \\ \hline
			Promedio                  & 6,72                \\ \hline
			Desviación estándar       & 0,083666            \\ \hline
			Valor $\pm$ Incertidumbre & 6,72 $\pm$ 0,083666 \\ \hline
		\end{tabular}
		\caption{Desviación estándar - Voltaje en AC [ Primera Medición ]}
	\end{center}
\end{table}

\begin{table}[H]
	\begin{center}
		% chktex-file 44
		\begin{tabular}{|c|c|}
			\multicolumn{2}{c}{Voltaje en AC}                 \\ \hline
			Promedio                  & 20,95                 \\ \hline
			Desviación estándar       & 0,618328              \\ \hline
			Valor $\pm$ Incertidumbre & 20,95  $\pm$ 0,618328 \\ \hline
		\end{tabular}
		\caption{Desviación estándar - Voltaje en AC [ Segunda Medición ]}
	\end{center}
\end{table}

\begin{table}[H]
	\begin{center}
		% chktex-file 44
		\begin{tabular}{|c|c|}
			\multicolumn{2}{c}{Voltaje en DC}               \\ \hline
			Promedio                  & 8,19                \\ \hline
			Desviación estándar       & 0,022361            \\ \hline
			Valor $\pm$ Incertidumbre & 8,19 $\pm$ 0,022361 \\ \hline
		\end{tabular}
		\caption{Desviación estándar - Voltaje en DC [ Primera Medición ]}
	\end{center}
\end{table}

\begin{table}[H]
	\begin{center}
		% chktex-file 44
		\begin{tabular}{|c|c|}
			\multicolumn{2}{c}{Voltaje en DC}                \\ \hline
			Promedio                  & 18,67                \\ \hline
			Desviación estándar       & 0,238747             \\ \hline
			Valor $\pm$ Incertidumbre & 18,67 $\pm$ 0,238747 \\ \hline
		\end{tabular}
		\caption{Desviación estándar - Voltaje en DC [ Segunda Medición ]}
	\end{center}
\end{table}

% ------------------------------------------------ |>
\subsection{Error en las mediciones de resistencias}

\begin{table}[H]
	\begin{center}
		% chktex-file 44
		\begin{tabular}{|c|c|c|}
			\multicolumn{3}{c}{Resistencias}                         \\ \hline
			Código        & Multímetro digital & Porcentaje de Error \\ \hline
			220 $\pm$ 5\% & 217,9 $\Omega$     & 0.9664 \%           \\ \hline
			150 $\pm$ 5\% & 149,1 $\Omega$     & 0.604  \%           \\ \hline
			470 $\pm$ 5\% & 465,0 $\Omega$     & 1.075  \%           \\ \hline
			330 $\pm$ 5\% & 324,3 $\Omega$     & 1.757  \%           \\ \hline
		\end{tabular}
	\end{center}
\end{table}

\subsection{Cálculos con la Ley de Ohm}

\begin{table}[H]
	\begin{center}
		% chktex-file 44
		\begin{tabular}{|c|c|c|c|c|c|c|c|c|c|}
			% Table Header
			\hline
			\multicolumn{3}{c}{Datos [Multimetro]}                 &
			\multicolumn{2}{c}{Valor [Voltios]}                    &
			\multicolumn{4}{c}{Resistencias [ Multimetro ]}          \\ \hline

			% 1ra fila
			\multicolumn{3}{c}{Voltaje en AC [ Primera Medicion ]} &
			\multicolumn{2}{c}{5,44}                               &
			217,90                                                 &
			149,10                                                 &
			465,00                                                 &
			324,30                                                   \\

			% 2da fila
			\multicolumn{3}{c}{Voltaje en AC [ Segunda Medición ]} &
			\multicolumn{2}{c}{17,17}                              &
			                                                       &
			                                                       &
			                                                       &
			\\

			% 3ra fila
			\multicolumn{3}{c}{Voltaje en DC [ Primera Medicion ]} &
			\multicolumn{2}{c}{6,88}                               &
			                                                       &
			                                                       &
			                                                       &
			\\

			% 4ta fila
			\multicolumn{3}{c}{Voltaje en DC [ Segunda Medicion ]} &
			\multicolumn{2}{c}{17,47}                              &
			                                                       &
			                                                       &
			                                                       &
			\\ \hline
		\end{tabular}
		\caption{Calculo de la corriente con la Ley de Ohm - Primera Parte}
		\label{tab:LeyDeOhmPrimeraParte}
	\end{center}
\end{table}

\begin{table}[H]
	\begin{center}
		% chktex-file 44
		\begin{tabular}{| p{2cm} | p{2cm} | p{2cm} | p{2cm} |}
			% Table Header
			\hline
			\multicolumn{4}{c}{Resultados [ I ]}      \\ \hline


			% 1ra fila
			1        & 2        & 3        & 4        \\ \hline

			% 2da fila
			0,024966 & 0,036486 & 0,011699 & 0,016775 \\ \hline

			% 3ra fila
			0,078798 & 0,115158 & 0,036925 & 0,052945 \\ \hline

			% 4ta fila
			0,031574 & 0,046144 & 0,014796 & 0,021215 \\ \hline

			% 5ta fila
			0,080174 & 0,117170 & 0,037570 & 0,053870 \\ \hline
		\end{tabular}
		\caption{Calculo de la corriente con la Ley de Ohm - Segunda Parte}
		\label{tab:LeyDeOhmSegundaParte}
	\end{center}
\end{table}

\begin{table}[H]
	\begin{center}
		\begin{tabular}{|cc|c|c|c|c|c|c|c|c|}
			% table Header
			\hline
			\multicolumn{2}{c}{Valor [ Voltios ]} &
			\multicolumn{4}{c}{Corriente [ I ]}   &
			\multicolumn{4}{c}{Resultado [ $\Omega$ ]}                                                            \\ \hline

			                                      &          & 1        & 2        & 3        & 4 & 1 & 2 & 3 & 4 \\ \hline

			\multicolumn{2}{c}{5,44}              & 0,024966 & 0,036486 & 0,011699 & 0,016775 &
			217,9                                 &
			149,1                                 &
			465                                   &
			324,3                                                                                                 \\

			\multicolumn{2}{c}{17,17}             & 0,078798 & 0,115158 & 0,036925 & 0,052945 &
			                                      &
			                                      &
			                                      &
			\\

			\multicolumn{2}{c}{6,88}              & 0,031574 & 0,046144 & 0,014796 & 0,021215 &
			                                      &
			                                      &
			                                      &
			\\

			\multicolumn{2}{c}{17,47}             & 0,080174 & 0,117170 & 0,037570 & 0,053870 &
			                                      &
			                                      &
			                                      &
			\\ \hline
		\end{tabular}
		\caption{Calculo de las resistencias con la Ley de Ohm}
		\label{tab:LeyDeOhmResistencias}
	\end{center}
\end{table}

% ------------------------------------------------ |>
\subsection{Respuesta a preguntas de la guía de laboratorio}

% -----------------------|>
\subsubsection{¿Cuáles son los pasos que debo hacer para realizar una medida
	con un multímetro?}

A pesar de ser el instrumento básico que cualquier persona relacionada con
la electrónica o la electricidad utiliza, la forma de utilizarlo cambia un poco
dependiendo de lo que queramos medir. Por ejemplo, para medir \textit{Voltaje},
la medición se realiza en paralelo, mientras que para medir \textit{Corriente},
la medición es en serie.

% -----------------------|>
\subsubsection{¿Cómo se conecta un multímetro para medir voltaje, corriente
	y resistencia?}

\begin{itemize}
	\item Voltaje: La medición de voltaje se realiza en paralelo, así que
	      sólo es necesario colocar la punta positiva del multímetro (roja) con
	      el punto positivo a medir, también hay que colocar la punta negativa
	      (negro) con el punto negativo a medir.~\cite{AcMaxMexico}

	\item Corriente: Lo primero a tener en cuenta es que la corriente no se
	      mide de la misma manera que el voltaje. La corriente se mide en serie.
	      Para hacer la medición en serie, hay que abrir nuestro circuito y
	      conectar el multímetro como si fuera un resistor más. La resistencia
	      del multímetro es muy pequeña y rara vez altera la exactitud de
	      nuestras mediciones.~\cite{AcMaxMexico}

	\item Resistencia: Para medir con un multímetro, sólo se colocan las puntas
	      en cada terminal de nuestra resistencia que queremos medir, sin importar
	      el color. Las puntas se colocan en el multímetro en el mismo lugar que
	      para medir voltaje y se selecciona la función con la letra griega omega
	      mayúscula $\Omega$.~\cite{AcMaxMexico}
\end{itemize}

% -----------------------|>
\subsubsection{¿Qué precauciones debo tener en cuenta con los instrumentos a
	la hora de realizar una medición?}

Según, (\cite{flukeCorporation}, Guía de Seguridad para Multímetros Digitales),
Antes de realizar una medida con el multímetro se debe someter a una inspección
visual. Compruebe que el multímetro, las sondas de prueba y los accesorios no
presentan daños físicos. Asegúrese de que todas las conexiones encajan
firmemente y de que no se aprecie metal al descubierto ni grietas en la carcasa.
No utilice nunca un multímetro ni sondas de prueba que presenten daños.

\vspace{.5cm}

Una vez finalizada la inspección visual, compruebe que el multímetro
funciona correctamente. Nunca se debe de dar por hecho.

% -----------------------|>
\subsubsection{¿Cuáles son las ventajas y desventajas del multímetro digital
	respecto al analógico?}

\begin{center}
	Ventajas
\end{center}

\begin{itemize}
	\item No muestra errores de lectura en la pantalla.
	\item En la mayoría de los multímetros, detección automática de
	      la polaridad.
	\item Detección automática del rango de medición.
	\item Mayor precisión del valor medido mostrado en la pantalla.
	\item Menos sensible a interferencias, ruidoso flujos magnéticos.
	\item Más económico de producir porque tiene menos componentes mecánicos.
\end{itemize}

\vspace{.5cm}

\begin{center}
	Desventajas
\end{center}

\begin{itemize}
	\item Con frecuencias altas existen mediciones incorrectas en corriente
	      alterna.
	\item Los picos de tensión pueden dañar el multímetro.
\end{itemize}

Fuente:~\cite{CesarCinjordiz}

% ----------------------------------------------------------------------|>
\section{Conclusiones}

Para concluir, podemos notar que en cuando se trate de resistencia toca
tener en cuenta los colores de esta, porque que toca observar bien los colores
para no medir mal cuál es valor de resistencia, por otro lado, la resistencia
siempre cuentan con un rango de error, Que al medir con un multímetro digital
podemos darnos cuenta de este. Por la parte del voltaje sabemos que el
multímetro digital es más exacto que el multímetro análogo. También podemos
notar que al medir el voltaje de corriente continua (dc) y el voltaje de
corriente alterna (ac), cuando se mide un voltaje bajo la diferencia entre
el multímetro digital y el análogo, la diferencia no es muy grande, pero al
medir un mayor voltaje la diferencia entre los dos multímetros se agranda.

\vspace{.5cm}

Por lo cual podemos decir que a mayor voltaje mayor es el rango de error
entre el multímetro digital y el análogo.

\newpage

\bibliography{./Bibliography/bibliography.bib}

\end{document}